%% ============================================================================
%% Fichier  : annexes/Ann_C_templates.tex
%% Titre    : Modèles et procédures
%% Version  : 1.0 -- novembre 2026
%% ============================================================================

\chapter{Modèles et Procédures}
\label{ann:templates}

Les modèles ci-dessous reprennent ceux des chapitres, sous une forme prête à l'emploi.

\section*{Fiche de mesure de sécurité}

\begin{tableaucours}{L{3.4cm}L{8.8cm}}{%
    \tableauHeader{Rubrique} &
    \tableauHeader{Contenu}
}
Identifiant et intitulé & M-xx, nom court \\
Objectif & Ce que la mesure empêche ou limite \\
Menace et risque couverts & Renvoi au registre des risques \\
Mise en œuvre & Ce qui est fait, par qui \\
Vérification & Test qui prouve que la mesure fonctionne, et sa fréquence \\
Preuve & Pièce conservée \\
\end{tableaucours}

\section*{Note d'autorisation d'un test}

\begin{tableaucours}{L{3.4cm}L{8.8cm}}{%
    \tableauHeader{Rubrique} &
    \tableauHeader{Contenu}
}
Signataires & Responsable du système, hébergeur ou fournisseur \\
Périmètre & Adresses, applications, comptes~; exclusions \\
Méthodes & Autorisées et interdites \\
Fenêtre & Dates et heures \\
Contacts & Urgence des deux côtés, procédure d'arrêt \\
Données & Traitement des données rencontrées et des résultats \\
\end{tableaucours}

\section*{Registre des risques}

\begin{tableaucours}{L{3.4cm}L{8.8cm}}{%
    \tableauHeader{Colonne} &
    \tableauHeader{Contenu}
}
Référence et scénario & Source, événement, conséquence \\
Cotation & Vraisemblance, impact, score, zone \\
Traitement & Option et mesures \\
Résiduel & Cotation après mesures vérifiées \\
Propriétaire & Personne qui répond du risque \\
Décision & Acceptation écrite, date \\
\end{tableaucours}

\section*{Procès-verbal de restauration}

\begin{tableaucours}{L{3.4cm}L{8.8cm}}{%
    \tableauHeader{Rubrique} &
    \tableauHeader{Contenu}
}
Système et source & Système restauré, copie utilisée \\
Durée & Heure de début et de fin, comparée au RTO \\
Données & Date des dernières données, comparée au RPO \\
Contrôles & Vérifications fonctionnelles \\
Anomalies et actions & Écarts et suites \\
Signatures & Participants \\
\end{tableaucours}

\section*{Fiche de chaîne de conservation}

Identification de la pièce~; saisie (lieu, date, personne, cadre)~;
acquisition (outil, version, empreintes)~; scellé~; mouvements (de qui, à qui,
motif, empreinte vérifiée)~; remise contre récépissé (Chapitre~\ref{chap:investigation}).
